\section{The independent evaluation}
\label{sec:evaluation}

The evaluation separates two questions: whether an edit \emph{repairs} false abstention (branch R), and whether
the proposed \emph{pathway and explanation} survive intervention tests (branch M). Each branch can select a
different edit and dose. Success in one does not imply success in the other.

\paragraph{Replay and screening on E.}
After the agent finished, the evaluation harness replayed every submission on all of E. The repair screen
required a utility gain over B0 with a one-sided 97.5\% lower bound above zero and at least 20 false-abstention
bases in DE or RC. The pathway screen required an effect of at least 0.05 in the declared direction, with a
one-sided 97.5\% bound beyond zero and at least 20 bases in the relevant endpoint population. If no submission
met either screen, S, H and T would have remained unread. E outcomes remained exploratory because the agent
had used that split to develop the edits.

\paragraph{Selection on S.}
The entry rules were then evaluated on S. For repair, candidate and comparator doses were chosen to maximise
utility, and the qualifying candidate with the highest utility was selected; candidate ties favoured fewer
edited weights. For the pathway branch, dose selection maximised the declared signed effect, subject to the
pre-registered limit on degenerate outputs. Candidate selection used the standardised effect, separately from
repair. Dose ties favoured the setting closest to no change.

The repair candidate was submission 2, a rank-1 edit to the layer-35 attention output projection, at dose 2.
Its utility on S was 0.429, compared with 0.402 for the next submission and 0.303 for B0. Dose 2 was the top of
the frozen grid; larger doses were not searched. The pathway candidate was submission 1, the two-head edit,
at dose 1.5, with the layer-33 head designated as its primary component. Selected repair-comparator doses
were $-8$ for logit bias (B1), 1 for steering (B4), and $\beta = 1$ for orthogonalisation (B5).

\paragraph{Freezing the confirmation.}
The selected candidates, doses, comparator settings and analysis code were frozen before any H or T model
outputs were generated or read. This freeze, F2, is recorded at commit \texttt{02bb2f6}, with design hash
\texttt{d02d9f9c\ldots}. The confirmation run checked that commit before accessing H or T.

\paragraph{Decision criteria.}
H and T were each evaluated once. The repair criterion on H required a positive lower bound for the utility
gain, damage no greater than 0.05, degenerate outputs no greater than 0.02, and counterfactual following of at
least 0.90. The damage and counterfactual thresholds applied to point estimates; uncertainty was reported
alongside them.

A better trade-off than a repair comparator could be established in either of two ways:
\textbf{O1}, higher recovery with an unsupported-answer increase smaller than 0.05; or
\textbf{O2}, fewer unsupported answers with a recovery loss smaller than 0.05. Superiority and non-inferiority
had to pass together for the relevant orientation. Holm correction covered all ten comparisons---two
orientations for each of B1--B5. Pathway tests formed a separate family (\cref{sec:explain}).

\paragraph{Uncertainty.}
We used a paired bootstrap clustered by base and stratified by task, with $B = 10{,}000$ draws.
Intervals are 95\% unless stated otherwise. The intervals describe uncertainty in individual estimates;
confirmatory decisions use the corrected tests. \Cref{app:methods} gives the inference rules.
